\documentclass[12pt]{article}
\usepackage[utf8]{inputenc}
\usepackage[T1]{fontenc}
\usepackage[margin=1in]{geometry}
\usepackage{setspace}
\usepackage{xcolor}
\usepackage[colorlinks=true, urlcolor=blue!60!black, linkcolor=blue!60!black]{hyperref}
\setlength{\parindent}{2em}
\setlength{\parskip}{1.25em}
\renewcommand{\baselinestretch}{1.0}
\begin{document}
\begin{center}
{\LARGE \textit{Chae Chan Ping v. United States}, 130 U.S. 581 (1889)\par}
\vspace{0.8em}
{\large Case Brief\par}
\end{center}
\vspace{1em}

\section{Memory Jogger}
In \textit{The Chinese Exclusion Case}, a unanimous Court upheld the Scott Act's bar on the return of Chinese laborers holding reentry certificates, holding that a later statute overrides an earlier conflicting treaty and that Congress's power to exclude noncitizens is an inherent incident of sovereignty whose exercise is conclusive on the courts (\href{https://supreme.justia.com/cases/federal/us/130/581/}{Justia}).

\section{Detailed Case Facts}
Chae Chan Ping, a subject of the Emperor of China and a laborer by occupation, resided lawfully in San Francisco, California from 1875 until June 2, 1887 (\href{https://www.law.cornell.edu/supremecourt/text/130/581}{Cornell LII}). Under the Burlingame Treaty of 1868, the United States and China had recognized ``the inherent and inalienable right of man to change his home and allegiance'' and mutual free migration. Following intense anti-Chinese political agitation on the Pacific Coast, the two nations concluded the Angell (Supplementary) Treaty of November 17, 1880, permitting the United States to ``regulate, limit, or suspend''---though not ``absolutely prohibit''---the immigration of Chinese laborers, while guaranteeing that Chinese laborers already in the United States could ``go and come of their own free will and accord.''

To execute the 1880 treaty, Congress enacted the Chinese Exclusion (Restriction) Act of May 6, 1882 (22 Stat.\ 58), as amended on July 5, 1884 (23 Stat.\ 115), suspending new Chinese labor immigration for ten years while directing customs collectors to issue reentry certificates to departing resident laborers as the ``only evidence permissible'' to reenter. Before departing San Francisco for China aboard the steamship \textit{Gaelic} on June 2, 1887, Chae Chan Ping duly obtained such a reentry certificate from the San Francisco Collector of Customs.

On September 7, 1888, Chae Chan Ping embarked on his return voyage from Hong Kong aboard the steamship \textit{Belgic}. While he was at sea, Congress passed the Scott Act of October 1, 1888 (25 Stat.\ 504), barring all departed Chinese laborers from returning and declaring every previously issued reentry certificate ``void and of no effect.'' When the \textit{Belgic} reached San Francisco on October 8, 1888, the Collector refused Chae Chan Ping permission to land solely on the basis of the Scott Act, and Captain Walker detained him aboard ship (\href{https://www.law.cornell.edu/supremecourt/text/130/581}{Cornell LII}).

\section{Procedural History}
Chae Chan Ping petitioned the U.S. Circuit Court for the Northern District of California for a writ of habeas corpus, arguing that the Scott Act violated existing U.S.--China treaties and unconstitutionally stripped him of a vested right under the Fifth Amendment. The Circuit Court denied the writ and remanded him to custody (36 F.\ 431 (1888)), and Chae Chan Ping appealed to the U.S. Supreme Court.

\section{Judicial Votes}
The Supreme Court unanimously \textbf{affirmed} the Circuit Court's denial of habeas corpus, \textbf{8--0}. The case was argued on March 28--29, 1889, and decided on May 13, 1889 (\href{https://supreme.justia.com/cases/federal/us/130/581/}{Justia}; \href{https://www.law.cornell.edu/supremecourt/text/130/581}{Cornell LII}):

\begin{itemize}
    \item \textbf{Opinion of the Court:} Justice Field delivered the opinion for a unanimous Court. The ninth seat was vacant after Justice Stanley Matthews died on March 22, 1889.
    \item \textbf{Concurring Opinions:} None.
    \item \textbf{Dissenting Opinions:} None.
\end{itemize}

(Four years later, however, Justice Field himself dissented in \textit{Fong Yue Ting v. United States}, 149 U.S.\ 698 (1893), arguing that while Congress enjoyed plenary power to \textit{exclude} noncitizens at the border, resident noncitizens already inside the United States could not be \textit{deported} without Fifth and Sixth Amendment protections (\href{https://supreme.justia.com/cases/federal/us/149/698/}{Justia}).)

\section{Holding}
\textbf{Questions Presented:}

\begin{enumerate}
  \item \textbf{Treaty Abrogation (Later-in-Time Rule):} Whether a federal statute that directly contravenes earlier treaties between the United States and a foreign nation is invalid or unenforceable in domestic courts.
  \item \textbf{Federal Exclusion Power and Judicial Review:} Whether Congress possesses constitutional authority to exclude noncitizens---including returning long-term residents holding statutory reentry certificates---and to what extent the Judiciary may review Congress's determination that exclusion is necessary.
\end{enumerate}

\textbf{Holding:} No on the first issue; yes on the second. In a unanimous 9--0 opinion by Justice Field, the Court affirmed the denial of habeas corpus and upheld the Scott Act of 1888 (\href{https://supreme.justia.com/cases/federal/us/130/581/}{Justia}). First, under the Supremacy Clause, federal statutes and treaties stand on equal constitutional footing, so a later-in-time Act of Congress abrogates prior conflicting treaty stipulations as domestic municipal law. Second, the power to exclude noncitizens is an inherent, inalienable incident of national sovereignty belonging exclusively to the federal government, and the legislative department's determination to exclude classes of noncitizens is conclusive upon the Judiciary.

\section{Analysis of Opinions}
\subsection{Opinion of the Court (Justice Field)}

\subsection{The Later-in-Time Rule and Nonjusticiability of Treaty Breaches}
Justice Field conceded candidly that the Scott Act of 1888 violated express stipulations of the 1868 and 1880 treaties with China. Nevertheless, he held the statute fully valid and enforceable domestically. Under Article VI's Supremacy Clause, federal statutes and treaties are both ``the supreme law of the land,'' with neither enjoying paramount domestic rank over the other. Relying on the \textit{Head Money Cases}, 112 U.S.\ 580 (1884), and \textit{Whitney v. Robertson}, 124 U.S.\ 190 (1888), Field reasoned that whenever a self-executing treaty and a federal statute conflict, ``the last expression of the sovereign will must control'' (\href{https://www.law.cornell.edu/supremecourt/text/130/581}{Cornell LII}). Whether Congress had sufficient cause to abrogate a treaty---and any diplomatic grievance China might assert---are political questions confided to the executive and legislative branches, not the courts.

\subsection{Inherent Sovereignty and the Birth of the Plenary Power Doctrine}
Turning to Congress's affirmative constitutional power to pass the Scott Act, Justice Field notably did not anchor immigration control solely in any single enumerated Article I clause (such as the Foreign Commerce or Naturalization Clauses). Instead, invoking Chief Justice Marshall's territorial-sovereignty reasoning in \textit{The Schooner Exchange v. M'Faddon}, 7 Cranch 116 (1812), and national-unity passages from \textit{Cohens v. Virginia}, 6 Wheat.\ 264 (1821), and \textit{Knox v. Lee}, 12 Wall.\ 457 (1871), Field derived the exclusion power from the nature of independent nationhood itself. For domestic matters, power is divided between States and the Union, ``but for national purposes, embracing our relations with foreign nations, we are but one people, one nation, one power'' (\href{https://www.law.cornell.edu/supremecourt/text/130/581}{Cornell LII}).

Because territorial sovereignty includes the absolute right to exclude foreigners as a matter of national self-preservation, Congress may exercise that power in peacetime no less than in war. Employing explicitly racialized rhetoric reflecting Pacific Coast nativism, Field asserted that if Congress deems ``the presence of foreigners of a different race in this country, who will not assimilate with us, to be dangerous to its peace and security,'' its judgment is ``conclusive upon the judiciary'' (\href{https://www.law.cornell.edu/supremecourt/text/130/581}{Cornell LII}).

\subsection{Inalienability and Rejection of Vested Reentry Rights}
Finally, the Court rejected Chae Chan Ping's claim that his statutory reentry certificate conferred an irrevocable vested property right. Sovereign powers are held in public trust and cannot be bargained away by treaty or statute. Distinguishing vested real-property or inheritance rights acquired under a treaty (\textit{Society for the Propagation of the Gospel v. New Haven}, 8 Wheat.\ 464 (1823)), Field characterized a statutory reentry certificate as merely a personal, revocable license ``held at the will of the government, revocable at any time, at its pleasure'' (\href{https://www.law.cornell.edu/supremecourt/text/130/581}{Cornell LII}).

\subsection{Subsequent Developments}

\begin{itemize}

\item \textbf{Expansion of the Plenary Power:} Immediately after 1889, the Court expanded \textit{Chae Chan Ping} into a sweeping edifice of immigration exceptionalism: \textit{Nishimura Ekiu v. United States}, 142 U.S.\ 651 (1892) (\href{https://supreme.justia.com/cases/federal/us/142/651/}{Justia}) (executive officers' exclusion decisions constitute due process for arriving noncitizens); \textit{Fong Yue Ting v. United States}, 149 U.S.\ 698 (1893) (\href{https://supreme.justia.com/cases/federal/us/149/698/}{Justia}) (extending plenary power from border exclusion to interior deportation); and \textit{United States ex rel.\ Knauff v. Shaughnessy}, 338 U.S.\ 537 (1950), and \textit{Shaughnessy v. United States ex rel.\ Mezei}, 345 U.S.\ 206 (1953) (upholding indefinite exclusion without a hearing on confidential security grounds).

\item \textbf{Limits:} Even when the Court recognized limited boundaries---such as forbidding criminal hard labor without a jury trial in \textit{Wong Wing v. United States}, 163 U.S.\ 228 (1896) (\href{https://supreme.justia.com/cases/federal/us/163/228/}{Justia}), or requiring a ``facially legitimate and bona fide reason'' when visa denials implicate U.S.\ citizens' First Amendment rights in \textit{Kleindienst v. Mandel}, 408 U.S.\ 753 (1972)---the core deference survived.

\item \textbf{Modern Reaffirmation:} In \textit{Trump v. Hawaii}, 585 U.S.\ 667 (2018), the Court applied deferential \textit{Mandel}-style rational-basis review to uphold Presidential Proclamation 9645 restricting entry from several predominantly Muslim nations.

\end{itemize}

\section{Examples: Future Applications}
\begin{enumerate}

\item \textbf{Same-Side (Statute Controls): Revoking a Visa-Waiver Treaty.} The United States ratifies a bilateral treaty guaranteeing citizens of Country X visa-free entry for business visits. Two years later, citing economic competition and security fears, Congress enacts a statute barring all Country X nationals from entering the United States and voiding all previously issued visas. A Country X national detained at an airport with a valid visa seeks habeas corpus. Under \textit{Chae Chan Ping}, the habeas petition fails. Congress's later-in-time statute abrogates the conflicting treaty as domestic law, the visa is a revocable license rather than a vested property right, and Congress's exclusion decision is largely conclusive upon the Judiciary.

\item \textbf{Same-Side (Later Statute Controls): A Tax Statute Overriding a Treaty.} A self-executing tax treaty between the United States and Country Y exempts Y's residents from U.S. tax on certain interest income. Congress later enacts a statute taxing that income for all foreign recipients, expressly including those exempt under earlier treaties. A Y resident sues for a refund, arguing that the treaty guarantees the exemption. Under the later-in-time rule that \textit{Chae Chan Ping} applied, a treaty and a federal statute stand on equal footing as the supreme law of the land, so ``the last expression of the sovereign will must control.'' The statute governs in U.S. courts; whether the United States has breached its international obligations is a matter for the political branches and diplomacy, not the courts (\href{https://supreme.justia.com/cases/federal/us/130/581/}{Justia}).

\item \textbf{Opposite-Side (Plenary Power Does Not Shield the Statute): Imprisonment at Hard Labor Without a Jury.} Suppose Congress enacts a statute authorizing immigration officers, without a jury trial or criminal indictment, to sentence noncitizens found unlawfully present in the interior to two years of imprisonment at hard labor prior to deportation. Under \textit{Wong Wing v. United States}, 163 U.S.\ 228 (1896), while Congress may detain or deport noncitizens administratively, imposing infamous criminal punishment inside the United States triggers full Fifth and Sixth Amendment criminal-procedure protections.

\item \textbf{Opposite-Side (No Deference to the Executive): A Proclamation Contrary to Statute.} Suppose the President issues an executive proclamation barring all asylum applications at the border in direct contradiction of an existing federal statute (8 U.S.C.\ \S\ 1158(a)(1)) granting a right to apply. \textit{Chae Chan Ping} vested plenary exclusion power in \textit{Congress} (the legislative department); when the President acts contrary to Congress's express statutory scheme, executive power is at its ``lowest ebb'' under \textit{Youngstown Sheet \& Tube Co.\ v. Sawyer}, and courts may enjoin ultra vires executive action without questioning Congress's sovereign authority.

\item \textbf{Fence-Sitter (Uncertain): A Statute Excluding Immigrants of One Religion.} Congress enacts a statute barring the admission of any immigrant who belongs to a named religion. A U.S. citizen whose spouse is denied a visa under the statute sues, invoking the Establishment and Equal Protection principles. \textit{Chae Chan Ping} treated Congress's exclusion judgments, even explicitly racial ones, as ``conclusive upon the judiciary,'' and \textit{Mandel} and \textit{Trump v. Hawaii} apply highly deferential review to entry decisions. But \textit{Trump v. Hawaii} emphasized that the proclamation there was neutral on its face, leaving open whether a facially religious exclusion would receive the same deference. Whether the plenary power tolerates an express religious exclusion has never been decided.

\end{enumerate}

\section{Critique}
\textbf{Extratextual Sovereignty vs.\ Enumerated Powers:} \textit{Chae Chan Ping} departs sharply from the foundational principle of \textit{McCulloch v. Maryland} that the federal government is one of enumerated powers. Deriving unenumerated ``inherent'' sovereignty from international law risks creating a constitutional black hole---what critics call a domain where ordinary constitutional rules have never applied---untethered from Article I limits and Bill of Rights constraints.

\textbf{Racial Animus and Unchecked Deference:} Field's overt reliance on racial unassimilability demonstrates the peril of declaring political-branch exclusion decisions ``conclusive upon the judiciary.'' Modern defenders argue the Court has cleansed the doctrine of racial justifications while retaining necessary foreign-affairs deference; critics counter that \textit{Trump v. Hawaii} and \textit{Mu\~{n}oz} show how easily plenary power still insulates discriminatory or arbitrary executive action.

\section{Key Quotations}
\begin{quote}
``By the constitution, laws made in pursuance thereof, and treaties made under the authority of the United States, are both declared to be the supreme law of the land, and no paramount authority is given to one over the other. ... In either case the last expression of the sovereign will must control.''
\end{quote}
\noindent --- Justice Field, 130 U.S.\ at 600 (\href{https://www.law.cornell.edu/supremecourt/text/130/581}{Cornell LII}).

\begin{quote}
``That the government of the United States, through the action of the legislative department, can exclude aliens from its territory is a proposition which we do not think open to controversy. Jurisdiction over its own territory to that extent is an incident of every independent nation. It is a part of its independence.''
\end{quote}
\noindent --- Justice Field, 130 U.S.\ at 603--604 (\href{https://www.law.cornell.edu/supremecourt/text/130/581}{Cornell LII}).

\begin{quote}
``For local interests the several states of the Union exist, but for national purposes, embracing our relations with foreign nations, we are but one people, one nation, one power. ... If, therefore, the government of the United States, through its legislative department, considers the presence of foreigners of a different race in this country, who will not assimilate with us, to be dangerous to its peace and security, their exclusion is not to be stayed because at the time there are no actual hostilities with the nation of which the foreigners are subjects. ... In both cases its determination is conclusive upon the judiciary.''
\end{quote}
\noindent --- Justice Field, 130 U.S.\ at 606 (\href{https://www.law.cornell.edu/supremecourt/text/130/581}{Cornell LII}).

\begin{quote}
``The power of exclusion of foreigners being an incident of sovereignty belonging to the government of the United States as a part of those sovereign powers delegated by the Constitution, the right to its exercise at any time when, in the judgment of the government, the interests of the country require it, cannot be granted away or restrained on behalf of anyone.''
\end{quote}
\noindent --- Justice Field, 130 U.S.\ at 609 (\href{https://supreme.justia.com/cases/federal/us/130/581/}{Justia}).


\end{document}
